\documentclass[10pt,a4paper]{amsart}
\usepackage[margin=19mm]{geometry}
\usepackage{amsmath,amssymb,mathtools}
\usepackage[scr=rsfs,cal=euler]{mathalfa}
\usepackage{xfrac}
\usepackage[unicode,hidelinks,bookmarks=false]{hyperref}
\newcommand{\ie}{i.e.\ }
\newcommand{\cf}{cf.\ }
\newcommand{\resp}{respectively\ }
\newcommand{\Subsection}{\subsection}
\providecommand{\nobreakdash}{\nobreak\hbox{-}}
\setlength{\emergencystretch}{2em}
\multlinegap0pt
\usepackage{algorithm}\usepackage{algpseudocode}
\usepackage{amssymb}
\usepackage{mathtools}
\newtheorem{theorem}{Theorem}
\newtheorem{definition}{Definition}
\newcommand{\C}{\mathbb C}
\newcommand{\diag}     {\mathop{\rm diag}\nolimits}
\newcommand {\mat}[1]{\left[\begin{array}{#1}}
\newcommand{\nnrm}[1]{{\left\vert\kern-0.25ex\left\vert\kern-0.25ex\left\vert #1
		\right\vert\kern-0.25ex\right\vert\kern-0.25ex\right\vert}}
\newcommand{\rank}     {\mathop{\rm rank}\nolimits}
\newcommand {\rix}{\end{array}\right]}
\title{Characterization of stability radii for robustly asymptotically stable dissipative Hamiltonian differential-algebraic systems}
\author{Peter Benner, Volker Mehrmann, Anshul Prajapati,, Punit Sharma}
\date{Source: arXiv:2605.13891v1 (CC BY 4.0)}
\begin{document}
\maketitle

\noindent\emph{Source and licence.} Characterization of stability radii for robustly asymptotically stable dissipative Hamiltonian differential-algebraic systems. Version arXiv:2605.13891v1 (CC BY 4.0), distributed by arXiv under the Creative Commons Attribution 4.0 International licence, http://creativecommons.org/licenses/by/4.0/, as recorded in the arXiv licence metadata for this exact version and shown on the abstract page https://arxiv.org/abs/2605.13891v1. Full licence notice: \texttt{licenses/arxiv-2605.13891-CC-BY-4.0.txt}.\par\smallskip

\noindent\textit{Editorial scope.} Counted unit: the article's theorem thm:dist-hiEJR (source0.tex lines 1121--1141). The article's complete original proof of it is delivered with it (source0.tex lines 1142--1285, one \texttt{proof} environment of 144 lines). No mathematics is written, repaired or re-derived: the statement and the proof are the article's own lines, in the article's own notation, and no retained line is substituted or re-flowed.  Retained uncounted as the source-local context the proof needs: lines 364--366; lines 497--506; lines 550--579; lines 752--773; lines 762--764; lines 776--783.  Nothing was omitted from any retained span, so the package carries no omission record.  The retained lines cite 1 works, whose entries are transcribed verbatim from the article's own bibliography source in the e-print and delivered with short editorial labels recorded in the item's reference mapping.  No retained line contains a figure environment, an image-inclusion command or drawing code, so no image is delivered and the package declares no asset dependency.  Editorial formatting only, in addition: an AMS \texttt{amsart} preamble that restates the article's own declarations and macros used by the retained lines, namely the environments \texttt{theorem}, \texttt{definition} on separate counters, together with the standard packages those lines need.  The retained environments print in this document as Theorem 1, Definition 1 in order of appearance, whereas the article prints the same items with the numbering of its own full document.  The complete native source of the article as distributed in the arXiv e-print for this exact version is bundled unchanged as \texttt{sources/200544/source0.tex}.
\begin{equation}\label{dhdae1}
E\dot x=(J-R) x,
\end{equation}

\begin{theorem}\label{thm:stabledh}
    Consider the  dH
    pencil $\lambda E-(J-R)$ associated with the dHDAE~\eqref{dhdae1}. Let $N(E)$ be a matrix with columns that span the right nullspace of $E$. Then the system is robustly asymptotically stable if and only if the following conditions hold:
    \begin{itemize}
        \item [a)] the pencil is regular;
        \item [b)] the index of the pencil is at most one and every principal submatrix of  $N(E)^*(J-R)N(E)$ is nonsingular.
        \item [c)] the maximal real part $\mu$ of a finite eigenvalue of $\lambda E-(J-R)$ (the spectral abscissa) is negative.
    \end{itemize}
    The boundary of the set of robustly asymptotically stable pencils are those pencils where one of the three conditions is violated.
\end{theorem}

\begin{definition}\label{def:distances}
	Consider a robustly  asymptotically stable dHDAE system of the form~\eqref{dhdae1}, and the following sets  of structured perturbations to the matrices $E, J$ and $R$,
	%
	\begin{align}
		\mathcal S_d(E,J,R)&:= \{(\Delta_E,\Delta_J,\Delta_R):\Delta_E\leq 0, E+\Delta_E\geq 0, \Delta_J^*=-\Delta_J, \Delta_R\leq 0, R+\Delta_R\geq 0 \}, \label{set:Sd}\\
		\mathcal S_i(E,J,R)&:= \{(\Delta_E,\Delta_J,\Delta_R): \Delta_E^*=\Delta_E,E+\Delta_E\geq 0, \Delta_J^*=-\Delta_J, \Delta_R^*=\Delta_R, R+\Delta_R\geq 0 \}.\label{set:Si}
	\end{align}
	%
	Let $\mathbb S\in \{\mathcal S_d,\mathcal S_i\}$. Then, we consider the  distances:
	%
	\begin{enumerate}
		\item The distance to singularity $d_{sing}^{\mathbb S}(E,J,R)$ under perturbations from  $\mathbb S\in \{\mathcal S_d,\mathcal S_i\}$, is defined by
		%
		\begin{equation}\label{def:dist-sing}
			d_{sing}^{\mathbb S}(E,J,R) := \inf\{\nnrm{(\Delta_E,\Delta_J,\Delta_R)}~:~(E+\Delta_E,J+\Delta_J-R-\Delta_R) \text{ is singular}\}.
		\end{equation}
		%
		\item The distance to higher index $d_{hi}^{\mathbb S}(E,J,R)$ under perturbations from  $\mathbb S\in \{\mathcal S_d,\mathcal S_i\}$, is defined by
		%
		\begin{equation}\label{def:dist-hi}
			d_{hi}^{\mathbb S}(E,J,R) := \inf\{\nnrm{(\Delta_E,\Delta_J,\Delta_R)}~:~(E+\Delta_E,J+\Delta_J-R-\Delta_R) \text{ has index }>1\}.
		\end{equation}
		%
		\item The distance to a system with eigenvalues on the imaginary axis $d_{im}^{\mathbb S}(E,J,R)$ under perturbations from the set $\mathbb S\in \{\mathcal S_d,\mathcal S_i\}$, is defined by
		%
		\begin{equation}\label{def:dist-im}
			d_{im}^{\mathbb S}(E,J,R) := \inf\{\nnrm{(\Delta_E,\Delta_J,\Delta_R)}~:~\Lambda(E+\Delta_E,J+\Delta_J-R-\Delta_R)\cap i\mathbb R \neq \emptyset \}.
		\end{equation}
	\end{enumerate}
\end{definition}

\begin{theorem}\label{map:herm}
	Let $x \in \C^n\setminus\{0\}$ and $y\in \C^{n}$. Then
	%
	\begin{itemize}
		\item [a)] there exists a Hermitian matrix $H \in \C^{n,n}$ such that $Hx = y$ if and only if $\Im{(x^*y)} = 0$ and we have
		\[
		\min\big\{\|H\|\; : H\in \C^{n,n},\; H^*=H, \; Hx=y\big\} = \frac{\|y\|}{\|x\|},
		\]
		and the minimum is attained by
		%
		\begin{equation}\label{def:hatH}
			\hat H_{(x,y)} := \frac{\|y\|}{\|x\|}\mat{cc}\frac{y}{\|y\|}&\frac{x}{\|x\|} \rix \mat{cc}\frac{y^*x}{\|x\|\|y\|}&1\\1&\frac{x^*y}{\|x\|\|y\|} \rix \mat{cc}\frac{y}{\|y\|}&\frac{x}{\|x\|} \rix^*
		\end{equation}
		%
		if $x$ and $y$ are linearly independent and by $\hat H_{(x,y)}:=\frac{yx^*}{x^*x}$, otherwise;
		\item [b)] there exists a skew-Hermitian matrix $S \in \C^{n,n}$ such that $Sx = y$ if and only if $\Re{(x^*y)} = 0$ and we have
		\[
		\min\big\{\|S\|\; :\; S ^*=-S, \; Sx=y\big\} = \frac{\|y\|}{\|x\|},
		\]
		and the minimum is attained by $\hat S:=-i\hat H_{(x,iy)}$, where $\hat{H}$ is defined in~\eqref{def:hatH}.
	\end{itemize}
\end{theorem}

		\begin{equation}
			\hat H_{(x,y)} := \frac{\|y\|}{\|x\|}\mat{cc}\frac{y}{\|y\|}&\frac{x}{\|x\|} \rix \mat{cc}\frac{y^*x}{\|x\|\|y\|}&1\\1&\frac{x^*y}{\|x\|\|y\|} \rix \mat{cc}\frac{y}{\|y\|}&\frac{x}{\|x\|} \rix^*
		\end{equation}

\begin{theorem}\label{map:def}
    Let $x,y \in \C^n \setminus \{0\}$. Then there exists a positive semidefinite Hermitian matrix $H$ such that $Hx=y$ if and only if $x^*y > 0$. If the latter condition is satisfied, then
    \[
    \min\left \{\|H\|:H \in \C^{n,n},H^*=H \geq 0,~Hx=y\right \}=\frac{{\|y\|}^2}{x^*y}
    \]
    and the minimum is attained for the rank one matrix
    $\tilde H=\frac{1}{x^*y}yy^*$.
\end{theorem}

\begin{theorem}\label{thm:dist-hiEJR}
    Consider a robustly asymptotically stable dHDAE system of the form~\eqref{dhdae1}. Then
    %
    \begin{enumerate}
        \item the distance to higher index $d_{hi}^{\mathcal S_d}(E,J,R)$ with respect to perturbations from the set $\mathcal S_{d}(E,J,R)$ is bounded via
        %
        \begin{equation}
            (d_{hi}^{\mathcal S_d}(E,J,R))^2 \leq \min_{1 \leq k \leq n} \inf_{x\neq 0} \left\{ \lambda_{n-k+1}(E)^2 + \frac{\|(N_k^*)^{\dagger} N_k^* J N_k x\|^2}{\|N_k x\|^2} + \frac{\|(N_k^*)^{\dagger} N_k^* R N_k x\|^4} {\big(x^* N_k^* R N_k x\big)^2} \right\},
        \end{equation}
        %
        where $N_k$ denotes a  matrix with  columns that span a basis of $\ker(E+\hat \Delta_E)$, $\hat \Delta_E$ is the optimal perturbation such that $\rank (E+\hat \Delta_E)=k$,
         and $\lambda_{n-k+1}(E)$ is the $(n-k+1)$th smallest eigenvalue of $E$.
        \item if $R>0$, the distance to higher index $d_{hi}^{\mathcal S_i}(E,J,R)$ with respect to perturbations from the set $\mathcal S_{i}(E,J,R)$ is bounded via
        %
        \begin{equation}
            (d_{hi}^{\mathcal S_i}(E,J,R))^2 \leq \min_{1\leq k \leq n}\inf_{x\neq 0} \left\{ \lambda_{n-k+1}(E)^2 + \frac{\|(N_k^*)^{\dagger} N_k^* J N_k x\|^2}{\|N_k x\|^2} + \frac{\|(N_k^*)^{\dagger} N_k^* R N_k x\|^2} {\|N_k x\|^2} \right\},
        \end{equation}
        %
        where $N_k$ denotes a  matrix with  columns that span a basis of $\ker(E+\hat \Delta_E)$, $\hat \Delta_E$ is the optimal perturbation such that $\rank (E+\hat \Delta_E)=k$, and $\lambda_{n-k+1}(E)$ is the $(n-k+1)$th smallest eigenvalue of $E$.
    \end{enumerate}
\end{theorem}

\begin{proof}
    Using Definition~\eqref{def:distances}, the distance to higher index under structured perturbations from the set $\mathbb S\in \{\mathcal S_d,\mathcal S_i\}$ is given by
    \[
    d_{hi}^{\mathbb S}(E,J,R) = \inf\Big\{\nnrm{(\Delta_E,\Delta_J,\Delta_R)} ~:~ (\Delta_E,\Delta_J,\Delta_R)\in \mathbb S,~ \text{ index of } (E+\Delta_E,J+\Delta_J-R-\Delta_R) > 1 \Big\}.
    \]
    In view of Theorem~\ref{thm:stabledh}, this can be equivalently written as
    %
    \begin{align*}
        &d_{hi}^{\mathbb S}(E,J,R) \\ &= \inf\Big\{ \nnrm{(\Delta_E,\Delta_J,\Delta_R)} ~:~ (\Delta_E,\Delta_J,\Delta_R)\in \mathbb S,~ \text{rank}\big(N(E+\Delta_E)^* (J+\Delta_J-R-\Delta_R) N(E+\Delta_E) \big) < n \Big\},
    \end{align*}
    %
    where $N(E+\Delta_E)$ denotes a basis matrix for the null space of $E+\Delta_E$.

    Observe that the constraint in the above formulation involves second-order perturbation terms through products of $\Delta_E$ with $\Delta_J$ and $\Delta_R$. To address this difficulty, we adopt a two-stage approach. First, we determine an optimal perturbation $\Delta_E$ such that $\dim \ker(E+\Delta_E) = k$. Then, for this fixed $\Delta_E$, we determine the minimal structured perturbations $(\Delta_J,\Delta_R)$ that satisfy the rank condition. Let $E = U \Sigma U^*$ be the spectral decomposition of $E$, with $\Sigma = \diag \{\lambda_1,\lambda_2,\ldots,\lambda_n\}$, where $\lambda_1 \ge \cdots \ge \lambda_n$. The optimal perturbation $\Delta_E$ achieving ${\rank}(E+\Delta_E)=n-k$ is given by
    \[
    \hat \Delta_E = - U \diag\{0,0,\ldots,0,\lambda_{n-k+1},\lambda_{n-k+2}, \ldots,\lambda_{n-1},\lambda_{n}\}U^*.% - \lambda_n x_n x_n^* - \lambda_{n-1} x_{n-1} x_{n-1}^* - \cdots - \lambda_{n-k+1} x_{n-k+1} x_{n-k+1}^*,
    \]
    %where $x_{n-k+1}, \ldots,x_n$ are the corresponding right eigenvectors of $E$.
    Let $N_k$ be a  matrix whose columns form a basis for $\ker(E+\hat \Delta_E)$.
    %With this choice of $\hat \Delta_E$,
    Then we obtain
    %
    \begin{align}\label{Sd:muintro}
        (d_{hi}^{\mathbb S}(E,J,R))^2 &\leq  \inf\Big\{ \lambda_{n-k+1}(E)^2 + \|\Delta_J\|^2 + \|\Delta_R\|^2 : (\hat \Delta_E,\Delta_j,\Delta_R)\in \mathbb S,\nonumber \\ & \hspace{3cm} \exists\, x\neq 0 \text{ such that } N_k^* (J+\Delta_J-R-\Delta_R) N_k x = 0 \Big\}
        \nonumber\\
        &=\lambda_{n-k+1}(E)^2 + \mu_k^{\mathbb S},
    \end{align}
    where
    %
    \begin{eqnarray}\label{hu_hi}
        \mu_k^{\mathbb S}:=\inf\Big\{ \|\Delta_J\|^2 + \|\Delta_R\|^2 : (\hat \Delta_E,\Delta_j,\Delta_R)\in \mathbb S,\exists\, x\neq 0 \text{ such that } N_k^* (J+\Delta_J-R-\Delta_R) N_k x = 0 \Big\}.
    \end{eqnarray}
    %
    Since $J+\Delta_J$ is skew-Hermitian and $R+\Delta_R$ is Hermitian positive semidefinite, the condition
    \[
    (J+\Delta_J - R - \Delta_R)y = 0
    \]
    holds if and only if
    \[
    (J+\Delta_J)y = 0 \quad \text{and} \quad (R+\Delta_R)y = 0.
    \]
    Therefore,
    %
    \begin{eqnarray}\label{eq:temp2}
        \mu_k^{\mathbb S} &= & \inf\Big\{  \|\Delta_J\|^2 + \|\Delta_R\|^2 : (\hat \Delta_E,\Delta_j,\Delta_R)\in \mathbb S, \exists\, x\neq 0 \text{ such that } \nonumber
        \\ && \hspace{4.5cm} N_k^* (J+\Delta_J) N_k x = 0,~ N_k^* (R+\Delta_R) N_k x = 0 \Big\}.
    \end{eqnarray}
    %
    In~\cite[Lemma 2.8]{MehMS16} it is shown that for a matrix $B$, the equation $B\Delta x = y$ holds if and only if $\Delta x = B^{\dagger}y$ and $BB^{\dagger}y = y$, where $B^{\dagger}$ denotes the Moore-Penrose pseudoinverse. Applying this to the above constraints in~\eqref{eq:temp2}, the relations
    \[
    N_k^* \Delta_J N_k x = - N_k^* J N_k x, \qquad N_k^* \Delta_R N_k x = - N_k^* R N_k x
    \]
    can be equivalently written as
    \[
    \Delta_J N_k x = - (N_k^*)^{\dagger} N_k^* J N_k x, \qquad \Delta_R N_k x = - (N_k^*)^{\dagger} N_k^* R N_k x,
    \]
    since the consistency conditions are trivially satisfied. Hence, we obtain
    %
    \begin{eqnarray}\label{Sd:high1}
        \mu_k^{\mathbb S} &=  \inf\Big\{  \|\Delta_J\|^2 + \|\Delta_R\|^2 : (\hat \Delta_E,\Delta_J,\Delta_R)\in \mathbb S,~\exists\, x\neq 0, \; \Delta_J N_k x = - (N_k^*)^{\dagger} N_k^* J N_k x, \;
        \nonumber \\ &\hspace {5cm}\Delta_R N_k x = - (N_k^*)^{\dagger} N_k^* R N_k x \Big\}.
    \end{eqnarray}
    %
    We now solve the associated minimal-norm mapping problems.
    First, let $\mathbb S=\mathcal S_d(E,J,R)$. By Theorem~\ref{map:herm} and~\cite[Theorem~2.3]{MehMS16}, there exist a skew-Hermitian matrix $\Delta_J$ and a Hermitian negative semidefinite matrix $\Delta_R\le 0$ satisfying the constraints in~\eqref{Sd:high1} if and only if
    \[
    x^* N_k^* J N_k x \in i\mathbb{R}
    \quad \text{and} \quad
    x^* N_k^* R N_k x > 0.
    \]
    These conditions are automatically satisfied due to the structural properties of the matrices $J$ and $R$.
    %
    Among all such admissible mappings, the minimal norms are given by
    \[
    \|\Delta_J\| = \frac{\|(N_k^*)^{\dagger} N_k^* J N_k x\|}{\|N_k x\|}, \qquad
    \|\Delta_R\| = \frac{\|(N_k^*)^{\dagger} N_k^* R N_k x\|^2}{x^* N_k^* R N_k x},
    \]
    and are attained by
    \[
    \Delta_J = i\,\hat H_{(N_k x,i(N_k^*)^{\dagger} N_k^* J N_k x)}, \qquad
    \Delta_R = -\frac{\big((N_k^*)^{\dagger} N_k^* R N_k x\big) \big((N_k^*)^{\dagger} N_k^* R N_k x\big)^*}{x^* N_k^* R N_k x},
    \]
    where $\hat H$ is defined in~\eqref{def:hatH}.
    %
    If $R\ge 0$ is singular and $x^* N_k^* R N_k x = 0$, then necessarily $(N_k^*)^{\dagger} N_k^* R N_k x = 0$. In this case, adopting the convention $0/0 = 0$, the choice $\Delta_R = 0$ satisfies both the mapping constraint in~\eqref{Sd:high1} and the minimal-norm condition. Moreover, the resulting perturbation satisfies $R + \Delta_R \ge 0$; see~\cite[Lemma~4.1]{MehMS16}.
    %
    Combining these observations in~\eqref{Sd:high1}, we obtain
    \begin{equation}\label{Sd:high11}
    \mu_k^{\mathcal S_d} =  \inf_{x \neq 0} \left\{
    \frac{\|(N_k^*)^{\dagger} N_k^* J N_k x\|^2}{\|N_k x\|^2} +
    \frac{\|(N_k^*)^{\dagger} N_k^* R N_k x\|^4}{\big(x^* N_k^* R N_k x\big)^2} \right\},
    \end{equation}
       and thus from~\eqref{Sd:muintro} and~\eqref{Sd:high11}, we have
    \[
    (d_{hi}^{\mathcal S_d}(E,J,R))^2 \le \min_{1 \le k \le n} \inf_{x \neq 0} \left\{ \lambda_{n-k+1}(E)^2 +
    \frac{\|(N_k^*)^{\dagger} N_k^* J N_k x\|^2}{\|N_k x\|^2} +
    \frac{\|(N_k^*)^{\dagger} N_k^* R N_k x\|^4}{\big(x^* N_k^* R N_k x\big)^2} \right\}.
    \]
    Next, consider $\mathbb S=\mathcal S_i(E,J,R)$. By Theorem~\ref{map:herm}, there exist a skew-Hermitian matrix $\Delta_J$ and a Hermitian matrix $\Delta_R$ satisfying the constraints in~\eqref{Sd:high1} if and only if
    \[
    x^* N_k^* J N_k x \in i\mathbb{R}
    \quad \text{and} \quad
    x^* N_k^* R N_k x \in \mathbb{R},
    \]
    which again holds trivially due to the structure of $J$ and $R$.
    %
    The corresponding minimal norms are
    \[
    \|\Delta_J\| = \frac{\|(N_k^*)^{\dagger} N_k^* J N_k x\|}{\|N_k x\|},
    \qquad
    \|\Delta_R\| = \frac{\|(N_k^*)^{\dagger} N_k^* R N_k x\|}{\|N_k x\|},
    \]
    and are attained by
    \[
    \Delta_J = -i\,\hat H_{(N_k x,\, -i (N_k^*)^{\dagger} N_k^* J N_k x)},
    \qquad
    \Delta_R = \hat H_{(N_k x,\,-(N_k^*)^{\dagger} N_k^* R N_k x)},
    \]
    where $\hat H$ is defined in~\eqref{def:hatH}. Note that $(R+\Delta_R)Nx=0$, hence on applying~\cite[Lemma 4.4]{MehMS16} we have $R+\Delta_R\geq 0$.
    %
    Applying these minimal-norm mappings to~\eqref{Sd:high1} yields
    \begin{equation}\label{Sd:high22}
    \mu_k^{\mathcal S_i} =  \inf_{x \neq 0} \left\{
    \frac{\|(N_k^*)^{\dagger} N_k^* J N_k x\|^2}{\|N_k x\|^2} +
    \frac{\|(N_k^*)^{\dagger} N_k^* R N_k x\|^2}{\|N_k x\|^2} \right\}.
    \end{equation}
    Thus in view of~\eqref{Sd:muintro} and~\eqref{Sd:high22}, we have
    \[
    (d_{hi}^{\mathcal S_i}(E,J,R))^2 \le \min_{1 \le k \le n} \inf_{x \neq 0} \left\{
    \lambda_{n-k+1}(E)^2 +
    \frac{\|(N_k^*)^{\dagger} N_k^* J N_k x\|^2}{\|N_k x\|^2} +
    \frac{\|(N_k^*)^{\dagger} N_k^* R N_k x\|^2}{\|N_k x\|^2} \right\}.
    \]
    This completes the proof.
    %
    % For $\mathbb S=\mathcal S_d(E,J,R)$, the minimal-norm structured mappings can be obtained from Theorem~\ref{map:herm} for optimal $\Delta_J^*=-\Delta_J$ and from~{Theorem~\ref{map:def}} for optimal $\Delta_R\leq 0$.d  The existense condition for these mappings are trivially satisfied. Using the minimal norm mappings we obtain
    % \[
    % (d_{hi}^{\mathcal S_d}(E,J,R))^2 \leq \min_{1\leq k \leq n}\inf_{x\neq 0} \left\{ \lambda_{n-k+1}(E)^2 + \frac{\|(N_k^*)^{\dagger} N_k^* J N_k x\|^2}{\|N_k x\|^2} + \frac{\|(N_k^*)^{\dagger} N_k^* R N_k x\|^4} {\big(x^* N_k^* R N_k x\big)^2} \right\}.
    % \]
    % In contrast, for $\mathbb S=\mathcal S_i(E,J,R)$, the corresponding minimal-norm mappings are characterized by Theorem~\ref{map:herm}, leading to
    % \[
    % (d_{hi}^{\mathcal S_i}(E,J,R))^2 \leq \min_{1\leq k \leq n}\inf_{x\neq 0} \left\{ \lambda_{n-k+1}(E)_n^2 + \frac{\|(N_k^*)^{\dagger} N_k^* J N_k x\|^2}{\|N_k x\|^2} + \frac{\|(N_k^*)^{\dagger} N_k^* R N_k x\|^2} {\|N_k x\|^2} \right\}.
    % \]
\end{proof}

\begin{thebibliography}{99}
\bibitem[MehMS1]{MehMS16} C.~Mehl, V.~Mehrmann, and P.~Sharma. \newblock Stability radii for linear {H}amiltonian systems with dissipation under structure-preserving perturbations. \newblock {\em SIAM Journal on Matrix Analysis and Applications}, 37(4):1625--1654, 2016.
\end{thebibliography}
\end{document}
